\subsection{MVDR 接收波束在 P0 中的全局最优性}

\noindent\textbf{定理。}考虑论文中的系统能耗最小化问题 $P0$。假设 $P0$ 存在全局最优解；对每个 UAV 和时隙，感知信道为上一命题所列的秩一模型，$\sigma^2>0$、$p_j(t)\ge0$、$\bm a_r(\theta_i(t))\ne\bm0$；感知发射、卸载发射和接收波束均属于 $\mathbb C^{N\times1}$，且能耗权重 $\omega_2\ge0$。感知发射波束允许乘以 $[0,1]$ 内的系数，并且不存在论文所列以外会阻止该缩放的波束约束。那么，$P0$ 至少有一个全局最优解，其每个 UAV 在每个时隙的接收波束均可取上一命题求得的单位范数 MVDR 向量。

\medskip
\noindent\textbf{证明。}任取 $P0$ 的一个可行解 $\tilde{\bm\zeta}$，保留其中所有时隙的轨迹、频谱匹配、CU 功率、卸载波束、卸载时长及计算频率。对任意 UAV $u_i$ 和时隙 $t$，记干扰加噪声矩阵为
\begin{equation}\label{eq:new-cov}
\mathbf R_i(t)=\sum_{j=1}^{J}\eta_{i,j}(t)p_j(t)\bm h_{c_j,u_i}^{H}(t)\bm h_{c_j,u_i}(t)+\sigma^2\mathbf I.
\end{equation}
由上一命题的推导，$\mathbf R_i(t)\succ0$，并且单位范数向量
\begin{equation}\label{eq:new-mvdr}
\bm g_i^{\rm MVDR}(t)=\frac{\mathbf R_i^{-1}(t)\bm a_r(\theta_i(t))}{\|\mathbf R_i^{-1}(t)\bm a_r(\theta_i(t))\|}
\end{equation}
在保持其他变量及感知发射波束不变时，使该 UAV 的感知 SINR 达到最大值。即使原发射目标增益为零，该向量仍为最大 SINR 的一个解。

记原可行解中的感知发射与接收波束分别为 $\tilde{\bm w}_i(t)$ 和 $\tilde{\bm g}_i(t)$，原感知 SINR 为
\begin{equation}\label{eq:new-old-sinr}
\tilde s_i(t)=\frac{\left|\tilde{\bm g}_i^H(t)\mathbf A(\theta_i(t))\tilde{\bm w}_i(t)\right|^2}{\displaystyle\sum_{j=1}^{J}\eta_{i,j}(t)p_j(t)\left|\tilde{\bm g}_i^H(t)\bm h_{c_j,u_i}^{H}(t)\right|^2+\sigma^2}.
\end{equation}
保持 $\tilde{\bm w}_i(t)$ 不变，先以 $\bm g_i^{\rm MVDR}(t)$ 替换原接收波束，记此时的感知 SINR 为 $\hat s_i(t)$。根据上述最大化结论，
\begin{equation}\label{eq:new-ineq}
\hat s_i(t)\ge\tilde s_i(t)\ge0.
\end{equation}
接着构造新的感知发射波束 $\bm w'_i(t)=c_i(t)\tilde{\bm w}_i(t)$，其中
\begin{equation}\label{eq:new-scale}
c_i(t)=\begin{cases}
\sqrt{\tilde s_i(t)/\hat s_i(t)},&\hat s_i(t)>0,\\
0,&\hat s_i(t)=0.
\end{cases}
\end{equation}
由式~\eqref{eq:new-ineq}，$0\le c_i(t)\le1$；当 $\hat s_i(t)=0$ 时，$\tilde s_i(t)=0$。因此新方案的感知 SINR $s'_i(t)$ 与感知发射功率满足
\begin{samepage}
\begin{align}\label{eq:new-same-sinr}
s'_i(t)&=c_i^2(t)\hat s_i(t)=\tilde s_i(t),\\
\|\bm w'_i(t)\|^2&=c_i^2(t)\|\tilde{\bm w}_i(t)\|^2\le\|\tilde{\bm w}_i(t)\|^2.\label{eq:new-power}
\end{align}
\end{samepage}
对所有 UAV 和时隙分别执行上述构造，所得方案记为 $\bm\zeta'$。下面验证其满足 $P0$ 的约束。

首先，接收波束满足单位范数约束；由式~\eqref{eq:new-power} 及原解可行性，感知发射功率仍不超过上限。由于式~\eqref{eq:new-same-sinr} 保持感知 SINR 不变，原有感知 SINR 阈值仍成立。论文模型的感知信息速率为
\begin{equation}\label{eq:new-rate-formula}
R_i^{\rm sen}(t)=\frac{\delta}{2\nu}\log_2\!\left(1+2\sigma_{\rm pre}^2\gamma^2B^3\nu\,s_i(t)\right),
\end{equation}
其中 $s_i(t)$ 为感知 SINR。由此新旧感知速率及数据量相等：
\begin{equation}\label{eq:new-same-rate}
R_i^{\rm sen\prime}(t)=\tilde R_i^{\rm sen}(t),\qquad
\bar D^{\rm sen}R_i^{\rm sen\prime}(t)=\bar D^{\rm sen}\tilde R_i^{\rm sen}(t).
\end{equation}
论文中的感知数据完整卸载和计算时延约束分别为
\begin{align}\label{eq:new-model-off}
D_{u_i}^{\rm off}(t) R_i^{\rm off}(t)&\ge\bar D^{\rm sen}R_i^{\rm sen}(t),\\
\bar D^{\rm sen}+D_{u_i}^{\rm off}(t)+\frac{C^{\rm sen}\bar D^{\rm sen}R_i^{\rm sen}(t)}{f_{u_i}(t)}&\le D_{\max}^{\rm sen}.\label{eq:new-model-delay}
\end{align}
由于卸载波束、卸载时长和计算频率均保持原值，感知任务卸载速率 $R_i^{\rm off}(t)$ 不变；结合原解的可行性，得
\begin{align}
\bar D^{\rm sen}R_i^{\rm sen\prime}(t)&=\bar D^{\rm sen}\tilde R_i^{\rm sen}(t)\le\tilde D_{u_i}^{\rm off}(t)\tilde R_i^{\rm off}(t),\label{eq:new-off}\\
\bar D^{\rm sen}+\tilde D_{u_i}^{\rm off}(t)+\frac{C^{\rm sen}\bar D^{\rm sen}R_i^{\rm sen\prime}(t)}{\tilde f_{u_i}(t)}&\le D_{\max}^{\rm sen}.\label{eq:new-delay}
\end{align}
故感知数据完整卸载与计算时延约束仍满足，且原模型中的 BS 感知计算能耗 $\kappa C^{\rm sen}\bar D^{\rm sen}R_i^{\rm sen}(t) f_{u_i}^2(t)$ 不变。

其次，对于共享 UAV $u_i$ 频谱的 CU $c_j$，感知阶段的 CU 卸载速率为
\begin{equation}\label{eq:new-cu-rate}
R_{i,j}^{(1)}(t)=B\log_2\!\left(1+\frac{p_j(t)|h_{c_j,\rm BS}(t)|^2}{|\bm h_{u_i,\rm BS}^H(t)\bm w_i(t)|^2+\sigma^2}\right).
\end{equation}
而新方案的感知发射干扰满足
\begin{equation}\label{eq:new-interference}
\left|\bm h_{u_i,\mathrm{BS}}^H(t)\bm w'_i(t)\right|^2=c_i^2(t)\left|\bm h_{u_i,\mathrm{BS}}^H(t)\tilde{\bm w}_i(t)\right|^2\le\left|\bm h_{u_i,\mathrm{BS}}^H(t)\tilde{\bm w}_i(t)\right|^2.
\end{equation}
由式~\eqref{eq:new-cu-rate} 和~\eqref{eq:new-interference}，第一阶段速率不降低，即 $R_{i,j}^{(1)\prime}(t)\ge\tilde R_{i,j}^{(1)}(t)$。第二、第三阶段速率及各阶段时长保持不变，而第一阶段在 CU 已卸载数据量中具有非负系数 $\eta_{i,j}(t)\bar D^{\rm sen}$，故 CU 成功卸载约束仍成立。其余变量和约束均未改变，包括跨时隙轨迹约束，故 $\bm\zeta'$ 是 $P0$ 的可行解。

最后比较目标值。令 $E_{P0}(\cdot)$ 表示系统长期平均加权能耗。BS 计算能耗、CU 能耗、UAV 卸载能耗与飞行能耗均不变，感知发射能耗为 $\bar D^{\rm sen}\|\bm w_i(t)\|^2$，故
\begin{align}\label{eq:new-objective}
E_{P0}(\bm\zeta')-E_{P0}(\tilde{\bm\zeta})
&=-\frac{\omega_2\bar D^{\rm sen}}{t_{\max}}\sum_{t=1}^{t_{\max}}\sum_{u_i\in\mathcal U}\bigl(1-c_i^2(t)\bigr)\left\|\tilde{\bm w}_i(t)\right\|^2\le0.
\end{align}
现令 $\bm\zeta^\star$ 为 $P0$ 的一个全局最优解，$E_{P0}^\star=E_{P0}(\bm\zeta^\star)$。对其执行上述构造，得到采用 MVDR 接收波束的可行解 $\bm\zeta^{\rm MVDR}$。由式~\eqref{eq:new-objective}，
\begin{equation}\label{eq:new-opt-upper}
E_{P0}(\bm\zeta^{\rm MVDR})\le E_{P0}(\bm\zeta^\star)=E_{P0}^\star.
\end{equation}
另一方面，由于 $\bm\zeta^{\rm MVDR}$ 对 $P0$ 可行，且 $E_{P0}^\star$ 是全局最小值，有
\begin{equation}\label{eq:new-opt-lower}
E_{P0}(\bm\zeta^{\rm MVDR})\ge E_{P0}^\star.
\end{equation}
结合上述两式可得 $E_{P0}(\bm\zeta^{\rm MVDR})=E_{P0}^\star$。因此，$\bm\zeta^{\rm MVDR}$ 是 $P0$ 的一个全局最优解，定理得证。\hfill$\square$

